\documentclass[11pt,a4paper,twoside,openright]{report}

\usepackage[top=25mm,bottom=25mm,right=25mm,left=30mm,head=12.5mm,foot=12.5mm]{geometry}
\let\openright=\cleardoublepage

\usepackage[a-2u]{pdfx}

\usepackage[
   backend=biber
%  ,style=iso-authoryear
  ,style=alphabetic
  ,citestyle=numeric
  ,sortlocale=cs_CZ
  ,bibencoding=UTF8
  %,block=ragged
]{biblatex}
\addbibresource{references.bib}

%% Přepneme na českou sazbu, fonty Latin Modern a kódování češtiny
\usepackage[czech]{babel}
\usepackage{lmodern}
\usepackage[T1]{fontenc}
\usepackage{textcomp}
\usepackage[utf8]{inputenc}

% Set fonts
\RequirePackage[osf]{mathpazo} % Palatino with oldstyle figures
\newcommand\liningnums[1]{\fontfamily{ppl}\selectfont#1}
\RequirePackage{eulervm}
\RequirePackage[scaled=.8819]{sourcecodepro} % Source Code Pro typeface for monospace

%%% Další užitečné balíčky (jsou součástí běžných distribucí LaTeXu)
\usepackage{amsmath}        % rozšíření pro sazbu matematiky
\usepackage{amsfonts}       % matematické fonty
\usepackage{amsthm}         % sazba vět, definic apod.
\usepackage{bm}             % tučné symboly (příkaz \bm)
\usepackage{graphicx}       % vkládání obrázků
\usepackage{fancyvrb}       % vylepšené prostředí pro strojové písmo
\usepackage{fancyhdr}       % prostředí pohodlnější nastavení hlavy a paty stránek
\usepackage{icomma}         % inteligetní čárka v matematickém módu
\usepackage{dcolumn}        % lepší zarovnání sloupců v tabulkách
\usepackage{booktabs}       % lepší vodorovné linky v tabulkách
\makeatletter
\@ifpackageloaded{xcolor}{
   \@ifpackagewith{xcolor}{usenames}{}{\PassOptionsToPackage{usenames}{xcolor}}
  }{\usepackage[usenames]{xcolor}} % barevná sazba
\makeatother
\usepackage{multicol}       % práce s více sloupci na stránce
\usepackage{caption}
\usepackage{enumitem}
\usepackage{lipsum}
\setlist[itemize]{noitemsep, topsep=0pt, partopsep=0pt}
\setlist[enumerate]{noitemsep, topsep=0pt, partopsep=0pt}
\setlist[description]{noitemsep, topsep=0pt, partopsep=0pt}
\usepackage{pdfpages}

\usepackage{tocloft}
\setlength\cftparskip{0pt}
\setlength\cftbeforechapskip{1.5ex}
\setlength\cftfigindent{0pt}
\setlength\cfttabindent{0pt}
\setlength\cftbeforeloftitleskip{0pt}
\setlength\cftbeforelottitleskip{0pt}
\setlength\cftbeforetoctitleskip{0pt}
\renewcommand{\cftlottitlefont}{\Huge\bfseries}
\renewcommand{\cftloftitlefont}{\Huge\bfseries}
\renewcommand{\cfttoctitlefont}{\Huge\bfseries}

%Barevná syntaxe kódu
\usepackage{listings}
\renewcommand{\lstlistlistingname}{Seznam kódů}% Listing -> Seznam kódů
\definecolor{bluekeywords}{rgb}{0,0,1}
\definecolor{greencomments}{rgb}{0,0.5,0}
\definecolor{redstrings}{rgb}{0.64,0.08,0.08}
\definecolor{xmlcomments}{rgb}{0.5,0.5,0.5}
\definecolor{types}{rgb}{0.17,0.57,0.68}
\lstset{language=[Sharp]C,
captionpos=b,
numbers=left,
frame=lines,
showspaces=false,
showtabs=false,
breaklines=true,
showstringspaces=false,
breakatwhitespace=true,
escapeinside={(*@}{@*)},
commentstyle=\color{greencomments},
morekeywords={partial, var, value, get, set},
keywordstyle=\color{bluekeywords},
stringstyle=\color{redstrings},
basicstyle=\ttfamily\small,
numbers=left, 
}
% vyznaceni odstavcu
\parindent=0pt
\parskip=11pt

% zakaz vdov a sirotku - jednoradkovych pocatku ci koncu odstavcu na prechodu mezi strankami
\clubpenalty=1000
\widowpenalty=1000
\displaywidowpenalty=1000

% nastaveni radkovani
\renewcommand{\baselinestretch}{1.20}

% nastavení hlavy a paty stránek
\fancyhf{}
%\renewcommand{\chaptermark}[1]{\markboth{#1}{}}
%\fancyhead[L]{\leftmark}
\fancyhead[C]{\AutorPrace, \Kruh}
\fancyfoot[R]{\thepage}
\fancyhead[L]{ÚPA semestrální práce}
\fancyhead[R]{\today}
%\renewcommand{\footrulewidth}{0pt}
%\fancypagestyle{plain}{%
%\fancyhf{} % clear all header and footer fields
%\fancyfoot[R]{\thepage}
%\renewcommand{\headrulewidth}{0pt}
%\renewcommand{\footrulewidth}{0.5pt}
%}
\fancypagestyle{plain}{
\fancyhead[C]{\AutorPrace, \Kruh}
\fancyhead[R]{\today}
\fancyfoot[R]{\thepage}
\fancyhead[L]{ÚPA semestrální práce}
}

% Tato makra přesvědčují mírně ošklivým trikem LaTeX, aby hlavičky kapitol
% sázel příčetněji a nevynechával nad nimi spoustu místa. Směle ignorujte.
\makeatletter
\def\@makechapterhead#1{
  {\parindent \z@ \raggedright 
   \Huge\bfseries \thechapter. #1
   \par\nobreak
   \vskip 20\p@
}}
\def\@makeschapterhead#1{
  {\parindent \z@ \raggedright 
   \Huge\bfseries #1
   \par\nobreak
   \vskip 20\p@
}}
\makeatother

% Trochu volnější nastavení dělení slov, než je default.
\lefthyphenmin=2
\righthyphenmin=2

% Zapne černé "slimáky" na koncích řádků, které přetekly, abychom si
% jich lépe všimli.
\overfullrule=1mm

%% Balíček hyperref, kterým jdou vyrábět klikací odkazy v PDF,
%% ale hlavně ho používáme k uložení metadat do PDF (včetně obsahu).
%% Většinu nastavítek přednastaví balíček pdfx.
\hypersetup{unicode}
\hypersetup{breaklinks=true}
\hypersetup{hidelinks}

%%% Prostředí pro sazbu kódu, případně vstupu/výstupu počítačových
%%% programů. (Vyžaduje balíček fancyvrb -- fancy verbatim.)

\DefineVerbatimEnvironment{code}{Verbatim}{fontsize=\small, frame=single}



\def\NazevPrace{Určení hodnoty Ludolfova čísla pomocí řady, kterou navrhl James Gregory}
\def\Kruh{159}
\def\AutorPrace{Vladislav Aulich}
\def\DatumOdevzdani{2021}

% Vedoucí práce: Jméno a příjmení s~tituly
\def\Vedouci{Ing. Jan Kohout}

% Fakulta a obor
\def\StudijniProgram{Fyzikální a výpočetní chemie}
\def\Fakulta{Fakulta chemicko-inženýrská}

\begin{document}
	
	%%% Titulní strana práce a další povinné informační strany

%%% Titulní strana práce

\pagestyle{empty}
\pagenumbering{gobble}
\hypersetup{pageanchor=false}

\begin{center}

\vspace{\stretch{8}}

\includegraphics[width=.8\textwidth]{img/logo}

\vspace{\stretch{8}}

{\Huge\bfseries\NazevPrace}

\vspace{8mm}
\mdseries{Semestrální práce}

\vspace{\stretch{8}}
\large
\begin{tabular}{rl}
Autor: & \AutorPrace \\
\noalign{\vspace{2mm}}
Kruh: & \Kruh\\
\noalign{\vspace{2mm}}
Studijní obor: & \StudijniProgram\\
\noalign{\vspace{2mm}}
Fakulta: & \Fakulta\\
\noalign{\vspace{2mm}}
Akademický rok: & 2021/2022\\
\noalign{\vspace{2mm}}
Předmět: & Úvod do programování a algoritmů \\
\noalign{\vspace{2mm}}
Vedoucí práce: & \Vedouci \\
\end{tabular}

\vspace{20mm}
Praha, \DatumOdevzdani
\end{center}

\pagenumbering{arabic}
	
	% Obsah
	\setcounter{tocdepth}{2}
	\tableofcontents
	
	\chapter{Teoretická část}
	\pagestyle{fancy}
	
	\section{Úvod}
	Ludolfovo číslo patří do množiny iracionálních čísel, nelze jej tedy vyjádřit jako podíl dvou celých čísel. 
	
	
	\section{Zadání úlohy}
	Určete hodnotu Ludolfova čísla pomocí řady, kterou navrhl James Gregory.
	
	\chapter{Implementace}
	\section{Postup řešení}
	\section{Zdrojový kód}
        \begin{lstlisting}[caption=Výsledný zdrojový kód]
        using System;
        using System.Collections.Generic;
        using System.Linq;
        using System.Text;
        using System.Threading.Tasks;
        
        namespace semestralka_konsole
        {
            class Program
            {
                /// <summary>
                /// Metoda, která pomocí řady Jamese Gregoryho
                /// stanový hodnotu Ludolfova čísla.
                /// </summary>
                /// <param name="n">Počet iterací (členů řady)</param>
                /// <returns>Hodnota pi datového typu decimal</returns>
                public static decimal VypocitejPi(uint n)
                {
                    decimal soucet = 0;
                    decimal pomocna = 0;
                    for (int i = 0; i < n; i++)
                    {
        
                        pomocna = 4m / (1 + (2 * i));
                        //Pro sudé koeficienty
                        if (i % 2 == 0)
                            soucet += pomocna;
                        //Pro liché koeficienty
                        else
                            soucet -= pomocna;
                    }
                    return soucet;
                }
                static void Main(string[] args)
                {
                    bool pokracovat = true;
                    uint n = 0;
                    do
                    {
                        Console.WriteLine("Zadejte počet iterací:");
                        try
                        {
                            n = uint.Parse(Console.ReadLine());
                            pokracovat = false;
                        }
                        catch (Exception e)
                        {
                            Console.WriteLine("Zadejte prosím celé kladné číslo menší než  4 294 967 295");
        
                        }
                    }
                    while (pokracovat);
                    //Provedení funkce pro uživatelem zadaný počet iterací
                    Console.WriteLine("Hodnota Ludolfova čísla pro {0} iterací je {1}", n, VypocitejPi(n));
                    Console.ReadKey();
                }
            }
        }
        \end{lstlisting}
	\section{Ukázka programu}
	\chapter*{Závěr}
	\pagestyle{empty}
	\addcontentsline{toc}{chapter}{Závěr}
	Při práci na projektu jsem se dozvěděl mnoho zajímavostí o historii čísla $\pi$. Zejména o jeho významu v matematice a snahách o stanovení jeho hodnoty. 
	
	Dále jsem rozšířil své znalosti o sázení zdrojového kódu v sázecím jazyce \LaTeX. Původně jsem kód vykresloval pomocí balíčku verbatim, u toho se mi však nepovedlo nastvit kód podle mých představ. Nakonec jsem tedy použil balíček \uv{listings}. U něj jsem narážel na problémy s kódem obsahujícím diakritiku.
	
	K tvorbě dokumentace jsem využíval webový nástroj Overleaf a jednotlivé verze jsem verzoval v repozitáři umístěném na serveru Github. Nové zkušenosti jsem získal i přizpůsobením si vývojového prostředí Visual Studio.
	
	Práci na projektu považuji za přínosnou, domnívám se, že jsem zadání splnil.

	%%% Seznam použité literatury
	\nocite{dokumentace}
	\nocite{medium}
	\nocite{citace}
	\nocite{Historie_pi}
	\nocite{wiki:pi}
	\nocite{zvyrazneni_kodu}
	\nocite{clisting}
	\printbibliography[title={Seznam použité literatury},heading={bibintoc}]
	
	%%% Seznam obrázků
	\openright
	\listoffigures
	\addcontentsline{toc}{chapter}{Seznam obrázků}
	
	%%% Seznam tabulek
	\clearpage
	\listoftables
	\lstlistoflistings
	\addcontentsline{toc}{chapter}{Seznam tabulek}
	
	%%% Přílohy k práci, existují-li. Každá příloha musí být alespoň jednou
	%%% odkazována z vlastního textu práce. Přílohy se číslují.
	
	%\part*{Přílohy}
	%\appendix
	
\end{document}
